% ============================================================
%  C: Como Programar -- 6ª Edição | Deitel & Deitel
%  Capítulo 2 -- Introdução à Programação em C
%  FAZENDO A DIFERENÇA (2.32 e 2.33)
% ============================================================
\documentclass[12pt,a4paper]{article}
\usepackage[utf8]{inputenc}
\usepackage[T1]{fontenc}
\usepackage[brazil]{babel}
\usepackage{geometry}
\geometry{top=2.5cm,bottom=2.5cm,left=3cm,right=3cm}
\usepackage{parskip}\setlength{\parskip}{6pt}\setlength{\parindent}{0pt}
\usepackage{xcolor,tcolorbox,enumitem,listings,fancyhdr,titlesec,hyperref,booktabs,array,amsmath}
\tcbuselibrary{skins,breakable}

\definecolor{deitelblue}{RGB}{0,70,127}
\definecolor{deitelgray}{RGB}{240,242,245}
\definecolor{deitelgreen}{RGB}{0,110,60}
\definecolor{deitellightblue}{RGB}{220,235,250}
\definecolor{deitelred}{RGB}{160,20,20}
\definecolor{deitellorange}{RGB}{200,90,0}
\definecolor{ecogreen}{RGB}{0,120,50}
\definecolor{ecolightgreen}{RGB}{210,240,220}

\newtcolorbox{boaprática}[1][]{colback=deitellightblue,colframe=deitelblue,
  fonttitle=\bfseries\small,title={Boa prática de programação},breakable,#1}
\newtcolorbox{respostabox}[1][]{colback=deitelgray,colframe=deitelblue!60,
  fonttitle=\bfseries\small\color{deitelblue},title={Resposta detalhada},breakable,#1}
\newtcolorbox{enunciadoBox}[1][]{colback=white,colframe=deitelblue,
  fonttitle=\bfseries,title={#1},breakable}
\newtcolorbox{saida}[1][]{colback=black!88,colframe=deitelblue,
  fontupper=\ttfamily\small\color{green!70},
  title={\small\color{white}Saída do programa},breakable,#1}
\newtcolorbox{pesquisa}[1][]{colback=ecolightgreen,colframe=ecogreen,
  fonttitle=\bfseries\small,title={Pesquisa externa realizada},breakable,#1}
\newtcolorbox{atencao}[1][]{colback=yellow!10,colframe=deitellorange,
  fonttitle=\bfseries\small,title={Atenção},breakable,#1}

\setlist[enumerate,1]{label=\textbf{\alph*)},leftmargin=2em}
\setlist[itemize]{leftmargin=2em}

\lstset{
  basicstyle=\ttfamily\small,backgroundcolor=\color{deitelgray},
  frame=single,framerule=0.4pt,rulecolor=\color{deitelblue!40},
  breaklines=true,showstringspaces=false,
  keywordstyle=\color{deitelblue}\bfseries,
  commentstyle=\color{deitelgreen}\itshape,
  stringstyle=\color{deitelred},
  language=C,numbers=left,numberstyle=\tiny\color{gray},
  xleftmargin=18pt,xrightmargin=5pt,stepnumber=1,
}

\pagestyle{fancy}\fancyhf{}
\fancyhead[L]{\small\color{deitelblue}\textbf{C: Como Programar -- 6ª ed. | Deitel \& Deitel}}
\fancyhead[R]{\small\color{deitelblue}Cap. 2 -- Fazendo a Diferença}
\fancyfoot[C]{\small\thepage}
\renewcommand{\headrulewidth}{0.4pt}

\titleformat{\section}[block]{\large\bfseries\color{deitelblue}}{\thesection.}{0.5em}{}[\titlerule]
\titleformat{\subsection}[block]{\normalsize\bfseries\color{deitelblue!80}}{}{}{}

\hypersetup{colorlinks=true,linkcolor=deitelblue,urlcolor=ecogreen}

\begin{document}

\begin{titlepage}
  \centering\vspace*{2cm}
  {\color{deitelblue}\rule{\linewidth}{2pt}}\\[0.4cm]
  {\huge\bfseries\color{deitelblue}C: Como Programar}\\[0.2cm]
  {\large\color{deitelblue}6ª Edição $\cdot$ Paul Deitel \& Harvey Deitel}\\[0.2cm]
  {\color{deitelblue}\rule{\linewidth}{2pt}}\\[1cm]
  {\LARGE\bfseries Capítulo 2}\\[0.3cm]
  {\Large\color{deitelblue!80}Introdução à Programação em C}\\[0.8cm]
  \begin{tcolorbox}[colback=ecolightgreen,colframe=ecogreen,width=0.85\linewidth]
    \centering{\large\bfseries\color{ecogreen}Fazendo a Diferença}\\[0.3cm]
    {\normalsize Exercícios 2.32 e 2.33\\[0.2cm]
    Calculadora de IMC\\
    Economia com transporte solidário}
  \end{tcolorbox}
  \vfill{\small Abordagem incremental Deitel}
\end{titlepage}

\tableofcontents\newpage

% ============================================================
\section{Exercício 2.32 -- Calculadora de IMC}
% ============================================================

\begin{enunciadoBox}[Enunciado 2.32]
Crie uma calculadora de IMC (Índice de Massa Corporal) que leia peso (kg) e
altura (m), calcule o IMC e exiba os intervalos de classificação do Ministério da
Saúde:
\begin{itemize}[noitemsep]
  \item Abaixo do peso: < 18,5
  \item Normal: 18,5 -- 24,9
  \item Acima do peso: 25 -- 29,9
  \item Obeso: $\geq 30$
\end{itemize}
\end{enunciadoBox}

\subsection*{Pesquisa: Fórmula de IMC}

\begin{pesquisa}
A fórmula do IMC, desenvolvida por Adolphe Quetelet no século XIX, é:
\[
\text{IMC} = \frac{\text{peso}}{\text{altura}^2}
\]
onde o peso está em kg e a altura em metros. O resultado é dado em kg/m². As
faixas de classificação são padronizadas pela OMS e adotadas pelo Ministério
da Saúde brasileiro.

\textbf{Limitações conhecidas:} o IMC não distingue entre massa muscular e gordura,
então atletas podem ser classificados erroneamente como ``acima do peso''. Para
crianças, idosos e gestantes, há tabelas específicas. Sempre é importante
consultar um profissional de saúde.
\end{pesquisa}

\subsection*{Solução em C}

\begin{respostabox}
\begin{lstlisting}
/* Ex2.32: imc.c
   Calculadora de IMC -- versao com inteiros */
#include <stdio.h>

int main( void )
{
    int peso;       /* peso em kg */
    int altura_cm;  /* altura em centimetros (para evitar frações) */
    int imc;        /* IMC inteiro (resultado truncado) */

    printf( "=== Calculadora de IMC ===\n" );
    printf( "Digite o peso (kg, numero inteiro): " );
    scanf( "%d", &peso );

    printf( "Digite a altura (cm, ex: 170 para 1.70m): " );
    scanf( "%d", &altura_cm );

    /* IMC = peso / altura^2, com altura em metros.
       Como usamos cm, multiplicamos por 10000 para
       compensar a conversao (1m = 100cm, e elevamos ao quadrado) */
    imc = ( peso * 10000 ) / ( altura_cm * altura_cm );

    printf( "\nSeu IMC (aproximado): %d\n", imc );

    printf( "\nVALORES DE IMC:\n" );
    printf( "  Abaixo do peso: menor que 18\n" );
    printf( "  Normal:         18 a 24\n" );
    printf( "  Acima do peso:  25 a 29\n" );
    printf( "  Obeso:          30 ou mais\n" );

    if ( imc < 18 ) {
        printf( "\nClassificacao: Abaixo do peso\n" );
    }
    if ( imc >= 18 && imc <= 24 ) {
        printf( "\nClassificacao: Normal\n" );
    }
    if ( imc >= 25 && imc <= 29 ) {
        printf( "\nClassificacao: Acima do peso\n" );
    }
    if ( imc >= 30 ) {
        printf( "\nClassificacao: Obeso\n" );
    }

    printf( "\nLembre: o IMC e apenas uma estimativa.\n" );
    printf( "Sempre consulte um profissional de saude.\n" );

    return 0;
}
\end{lstlisting}

\textbf{Exemplo de execução:}
\begin{saida}
Digite o peso (kg, numero inteiro): 70\\
Digite a altura (cm, ex: 170 para 1.70m): 175\\
\\
Seu IMC (aproximado): 22\\
\\
VALORES DE IMC:\\
\textvisiblespace\textvisiblespace Abaixo do peso: menor que 18\\
\textvisiblespace\textvisiblespace Normal:         18 a 24\\
\textvisiblespace\textvisiblespace Acima do peso:  25 a 29\\
\textvisiblespace\textvisiblespace Obeso:          30 ou mais\\
\\
Classificacao: Normal
\end{saida}
\end{respostabox}

\subsection*{Análise didática}

\begin{atencao}
\textbf{Por que usamos altura em cm e o truque ``×10000''?}

Como o Cap.~2 só apresenta o tipo \texttt{int}, não podemos usar valores fracionários
diretamente. O truque é trabalhar com altura em centímetros (inteiro) e ajustar a
fórmula:

\[
\text{IMC} = \frac{\text{peso (kg)}}{\left(\frac{\text{altura}_\text{cm}}{100}\right)^2} = \frac{\text{peso} \times 10000}{\text{altura}_\text{cm}^2}
\]

Para peso=70 kg e altura=175 cm:
\[
\text{IMC} = \frac{70 \times 10000}{175^2} = \frac{700000}{30625} = 22{,}86 \approx 22 \text{ (truncado)}
\]

Esta solução demonstra criatividade dentro das limitações de cada capítulo --
abordagem fundamental da progressão Deitel.

\textbf{No Capítulo~3}, com \texttt{double}, faremos:
\begin{lstlisting}
double peso, altura, imc;
imc = peso / ( altura * altura );  /* simples e direto */
\end{lstlisting}
\end{atencao}

\subsection*{Por que isso é ``fazer a diferença''?}

\begin{tcolorbox}[colback=ecolightgreen,colframe=ecogreen,
  title={\bfseries Tecnologia a serviço da saúde}]

A obesidade é considerada uma epidemia global pela OMS. Calculadoras de IMC,
embora simples, são ferramentas educativas valiosas que:

\begin{itemize}[noitemsep]
  \item Promovem \textbf{consciência sobre saúde} entre populações.
  \item Permitem \textbf{rastreamento populacional} de tendências (epidemiologia).
  \item Servem como \textbf{ponto de partida} para conversas com profissionais.
  \item Fazem parte de aplicativos como MyFitnessPal, Strava, FitBit -- usados
  por bilhões.
\end{itemize}

Você acaba de implementar, com apenas \texttt{int} e \texttt{if}, o núcleo de
uma feature presente em incontáveis aplicações de saúde.
\end{tcolorbox}

\newpage

% ============================================================
\section{Exercício 2.33 -- Economia com Transporte Solidário}
% ============================================================

\begin{enunciadoBox}[Enunciado 2.33]
Pesquise sobre transporte solidário (\textit{carpooling}). Crie um programa que
calcule o gasto diário com automóvel próprio, considerando:
\begin{itemize}[noitemsep]
  \item Total de km dirigidos por dia
  \item Custo por litro de combustível
  \item Consumo médio (km/L)
  \item Estacionamento diário
  \item Pedágios diários
\end{itemize}
Exiba o total do gasto diário.
\end{enunciadoBox}

\subsection*{Pesquisa: Transporte Solidário}

\begin{pesquisa}
O transporte solidário (\textit{carpooling}) é a prática de compartilhar veículo
entre pessoas que fazem trajetos similares, dividindo custos e reduzindo emissões.
Benefícios incluem:

\begin{itemize}[noitemsep]
  \item \textbf{Econômico:} redução de até 50\% nos gastos diários ao dividir
  custos com um colega.
  \item \textbf{Ambiental:} cada carro a menos nas ruas reduz $\sim 4$ kg de
  CO\textsubscript{2} por dia (para 30 km típicos).
  \item \textbf{Social:} oportunidade de criar laços com colegas/vizinhos.
  \item \textbf{Trânsito:} menos congestionamentos para todos.
\end{itemize}

Plataformas no Brasil: BlaBlaCar, Uber Pool (já descontinuado), apps corporativos
de carona, grupos em redes sociais.
\end{pesquisa}

\subsection*{Solução em C}

\begin{respostabox}
\begin{lstlisting}
/* Ex2.33: transporte_solidario.c
   Calcula gasto diario com automovel proprio */
#include <stdio.h>

int main( void )
{
    int km_dia;         /* km dirigidos por dia */
    int custo_litro;    /* em centavos (ex: 580 = R$ 5,80) */
    int km_por_litro;   /* consumo */
    int estacionamento; /* em centavos */
    int pedagios;       /* em centavos */
    int gasto_combustivel;
    int gasto_total;
    int gasto_dividido;  /* se dividir com mais 1 pessoa */

    printf( "=== Calculadora de Gastos com Transporte ===\n\n" );

    printf( "Km dirigidos por dia: " );
    scanf( "%d", &km_dia );

    printf( "Custo do litro de combustivel (em centavos, ex: 580): " );
    scanf( "%d", &custo_litro );

    printf( "Consumo medio (km/litro): " );
    scanf( "%d", &km_por_litro );

    printf( "Estacionamento diario (centavos): " );
    scanf( "%d", &estacionamento );

    printf( "Pedagios diarios (centavos): " );
    scanf( "%d", &pedagios );

    /* Calcula gastos */
    gasto_combustivel = ( km_dia * custo_litro ) / km_por_litro;
    gasto_total = gasto_combustivel + estacionamento + pedagios;
    gasto_dividido = gasto_total / 2;

    printf( "\n=== RELATORIO DIARIO ===\n" );
    printf( "Combustivel: %d centavos\n", gasto_combustivel );
    printf( "Estaciona.:  %d centavos\n", estacionamento );
    printf( "Pedagios:    %d centavos\n", pedagios );
    printf( "-------------------------\n" );
    printf( "TOTAL DIARIO:        %d centavos\n", gasto_total );
    printf( "Por mes (22 dias):   %d centavos\n", gasto_total * 22 );

    printf( "\n=== COM TRANSPORTE SOLIDARIO (2 pessoas) ===\n" );
    printf( "Gasto pessoal/dia:   %d centavos\n", gasto_dividido );
    printf( "ECONOMIA por dia:    %d centavos\n",
            gasto_total - gasto_dividido );
    printf( "ECONOMIA por mes:    %d centavos\n",
            ( gasto_total - gasto_dividido ) * 22 );

    return 0;
}
\end{lstlisting}

\textbf{Exemplo de execução:}
\begin{saida}
Km dirigidos por dia: 40\\
Custo do litro de combustivel (em centavos, ex: 580): 580\\
Consumo medio (km/litro): 10\\
Estacionamento diario (centavos): 1500\\
Pedagios diarios (centavos): 500\\
\\
=== RELATORIO DIARIO ===\\
Combustivel: 2320 centavos\\
Estaciona.:  1500 centavos\\
Pedagios:    500 centavos\\
-------------------------\\
TOTAL DIARIO:        4320 centavos\\
Por mes (22 dias):   95040 centavos\\
\\
=== COM TRANSPORTE SOLIDARIO (2 pessoas) ===\\
Gasto pessoal/dia:   2160 centavos\\
ECONOMIA por dia:    2160 centavos\\
ECONOMIA por mes:    47520 centavos
\end{saida}

(R\$ 475,20 de economia mensal -- com apenas mais uma pessoa!)
\end{respostabox}

\subsection*{Análise didática}

\begin{atencao}
\textbf{Por que tudo em centavos?}

Mesma estratégia do Ex.\ 2.32: como \texttt{int} não suporta decimais, usamos a
menor unidade monetária (centavos). Para exibir como reais, basta lembrar que
4320 centavos = R\$ 43,20. No Cap.~3, com \texttt{double}, poderemos trabalhar
com vírgulas diretamente.

\textbf{Lições importantes:}
\begin{itemize}[noitemsep]
  \item \textbf{Modelagem do problema:} 5 entradas, 3 cálculos intermediários, 1 saída final.
  \item \textbf{Reuso de cálculos:} \texttt{gasto\_total} é usado várias vezes.
  \item \textbf{Análise comparativa:} mostrar economia explícita é mais
  persuasivo que apenas o gasto.
\end{itemize}
\end{atencao}

\subsection*{Por que isso é ``fazer a diferença''?}

\begin{tcolorbox}[colback=ecolightgreen,colframe=ecogreen,
  title={\bfseries Pequenas economias, grande impacto}]

\textbf{Impacto econômico individual:} R\$ 475 de economia mensal correspondem
a \textbf{R\$ 5\,700 por ano}. Em 5 anos: R\$ 28\,500 -- o suficiente para um
carro popular usado!

\textbf{Impacto ambiental coletivo:} se 10\% dos motoristas brasileiros adotassem
carona compartilhada uma vez por semana:
\begin{itemize}[noitemsep]
  \item $\sim 4$ milhões de carros a menos nas ruas por dia
  \item Redução de $\sim 16$ milhões de kg de CO\textsubscript{2} por dia
  \item Redução significativa de congestionamentos em grandes cidades
\end{itemize}

\textbf{Programação como ferramenta de tomada de decisão:} ao quantificar custos
e economias, ferramentas como esta ajudam pessoas a tomar decisões \textit{informadas
e baseadas em dados}. Não dizem o que fazer -- mostram as consequências
financeiras e ambientais de cada escolha.
\end{tcolorbox}

\newpage

% ============================================================
\section*{Reflexão Final -- Capítulo 2}

\begin{tcolorbox}[colback=deitellightblue,colframe=deitelblue,
  title={\bfseries Os primeiros programas que importam}]

Pode parecer surreal: estamos no Capítulo~2, e já escrevemos programas que:

\begin{itemize}[noitemsep]
  \item Calculam estatísticas (Ex.\ 2.16, 2.19, 2.23)
  \item Identificam padrões matemáticos (Ex.\ 2.24, 2.26)
  \item Geram arte ASCII (Ex.\ 2.21, 2.22, 2.25, 2.27)
  \item Manipulam dígitos (Ex.\ 2.29, 2.30)
  \item Avaliam saúde pessoal (Ex.\ 2.32)
  \item Estimam economias econômicas e impacto ambiental (Ex.\ 2.33)
\end{itemize}

\textbf{Com apenas:} \texttt{int}, \texttt{printf}, \texttt{scanf}, operadores
aritméticos, \texttt{if}. Cinco ingredientes -- e já estamos resolvendo problemas
reais.

\textbf{Imagine o que faremos quando adicionarmos:}
\begin{itemize}[noitemsep]
  \item Cap.~3: \texttt{while}, \texttt{double}, atribuição composta
  \item Cap.~4: \texttt{for}, \texttt{switch}, operadores lógicos
  \item Cap.~5: funções, recursão, números aleatórios
  \item Cap.~6: arrays e ordenação
  \item Cap.~7: ponteiros (o poder real de C)
  \item Cap.~8 em diante: strings, estruturas, arquivos\ldots
\end{itemize}

A jornada está apenas começando. \textbf{Cada capítulo amplia o que você pode
expressar} -- e tudo se constrói sobre o que aprendemos aqui.

\end{tcolorbox}

\begin{boaprática}
  \textbf{Dica de aprendizado:} ao concluir cada capítulo, releia os exercícios
  ``Fazendo a Diferença'' e tente \textit{melhorá-los} com o que você acabou de
  aprender. Por exemplo, no Cap.~3, refaça a calculadora de IMC usando
  \texttt{double} -- ela ficará muito mais precisa e elegante. Essa prática de
  \textit{revisitar} solidifica o aprendizado e mostra a progressão.
\end{boaprática}

\end{document}
